% Contribuições: subseção da Introdução, como no template da disciplina. Só o que o trabalho entrega; os
% resultados ficam no último parágrafo da Introdução e na Discussão.

\subsection{Contribuições}\label{sec:contribuições}

\begin{itemize}
\item Uma transformação dos eventos da partida em uma tabela de ações, com colunas de contexto (zona do campo,
recebedor do passe, fase de posse e placar), e em um grafo, em que cada ação é uma aresta do jogador para a
partida (Seção~\ref{sec:dados}).
\item Um conjunto de \nQuestions{} perguntas sobre a final da Copa do Mundo de 2022, com gabarito
(Seção~\ref{sec:perguntas}).
O código, as perguntas e o gabarito são públicos (Seção~\ref{sec:setup}).
\item Uma comparação de \nArms{} abordagens de aprendizado em contexto sob as mesmas perguntas, o mesmo prompt e
a mesma saída, em \nModels{} modelos de tamanhos diferentes (Seção~\ref{sec:experimentos}).
\end{itemize}

O restante do texto está organizado da seguinte forma: a Seção~\ref{sec:trabalhos} situa o trabalho em relação
à literatura; a Seção~\ref{sec:metodologia} descreve os dados, as abordagens comparadas, as perguntas e a
conferência das respostas; a Seção~\ref{sec:experimentos} apresenta os resultados e a discussão; e a
Seção~\ref{sec:conclusão} conclui o trabalho.
